% SPDX-License-Identifier: AGPL-3.0-or-later
% © Concepts 1996–2026 Miroslav Šotek. All rights reserved.
% © Code 2020–2026 Miroslav Šotek. All rights reserved.
% ORCID: 0009-0009-3560-0851 | Contact: www.anulum.li | protoscience@anulum.li
% Paper 1 of 2 — the operational-protocol null result.
% Companion: papers/submissions/02_grid_modal_growth_regime_map/.
% Every number binds to the hash-sealed artefacts under examples/real_data/.

\documentclass[11pt]{article}

\usepackage[T1]{fontenc}
\usepackage[utf8]{inputenc}
\usepackage{lmodern}
\usepackage{microtype}
\usepackage{amsmath}
\usepackage{booktabs}
\usepackage[margin=1.05in]{geometry}
\usepackage[hidelinks]{hyperref}
\usepackage{url}
\urlstyle{same}
\emergencystretch=2em

\newcommand{\code}[1]{\texttt{#1}}

\title{Generic early-warning signals are at chance under a matched
false-alarm protocol across brain, heart, grid, climate and molecular
data}
\author{Miroslav \v{S}otek\\
  Anulum Institute\\
  ORCID 0009-0009-3560-0851 \quad \href{mailto:protoscience@anulum.li}{protoscience@anulum.li}}
\date{August 2026}

\begin{document}
\maketitle

\begin{abstract}
Generic early-warning signals (EWS) --- the rising variance and lag-one
autocorrelation of critical slowing down, and related synchronisation
and ordinal-entropy indicators --- are widely reported to precede abrupt
transitions in the brain, the heart, power systems and the climate. Most
of that evidence rests on a retrospective, per-record test: a rising
trend measured over one pre-transition window and compared to surrogates
of the same record. We evaluate the operational question instead: at a
fixed false-alarm budget, does an early-warning detector fire on
transitions more often than on no-transition controls? One
domain-adaptable detector suite and one matched-false-alarm harness are
applied unchanged, through per-domain adapters, to four independent
labelled corpora --- scalp-EEG seizures, cardiac atrial-fibrillation
onsets, power-grid growing oscillations, and palaeoclimate abrupt
transitions --- and every result is scored with a label-permutation
significance test. No detector reaches significance in any domain (every
best-member $p \ge 0.05$); the sparse detection observed is what the
matched false-alarm rate produces by chance. The canonical literature
detector --- the Dakos et al.\ 2008 AR(1)-Kendall-$\tau$ trend --- run
through the same protocol on the same segments fares no better: on its
own palaeoclimate records it leads 0 of 6 transitions ($p = 1.00$), and
its strongest showing anywhere is 3 of 6 on scalp EEG ($p = 0.067$).
Extending the protocol to a molecular fifth modality ---
dynamical-network-biomarker (DNB) indices --- reaches the same
conclusion under modality-appropriate nulls: the single-cell transition
index of Mojtahedi et al.\ 2016 clears a matched operating point on 1
of 3 lineages ($p = 0.27$), and the bulk GSE2565 index does not beat a
surrogate null granted the same module-selection freedom as the
analysis ($p = 0.39$). Because the corpora are small (3--12 transitions
per domain), these nulls bound the demonstrated skill rather than
proving early warning impossible; what they establish is that the
operational bar is materially stricter than the retrospective
per-record test on which most published EWS evidence rests. Every alarm
and non-detection is sealed in content-addressed, byte-reproducible
evidence records; a companion manuscript shows a domain-specific
detector clearing the identical bar where a deterministic instability
signature exists.
\end{abstract}

\section{Background and question}\label{sec:background}

The theory of critical transitions predicts that a dynamical system
approaching a bifurcation recovers more slowly from perturbations, which
shows up as a rising variance and a rising lag-one autocorrelation of an
observable --- critical slowing down (Scheffer et al.\ 2009). The same
idea motivates rising-synchronisation and ordinal-transition-entropy
indicators. These generic early-warning signals have been reported
before epileptic seizures, cardiac arrhythmias, power-system
instability, ecological collapse and abrupt climate change.

Almost all of that evidence answers a retrospective, per-record
question: \emph{within one pre-transition window, is there a
statistically significant rising trend of the indicator, relative to
surrogate time series of that same record?} (e.g.\ the Kendall-$\tau$
trend test of Dakos et al.\ 2008). This question is valuable, but it is
not the question an operator or clinician faces. The operational
question is: \emph{if I set an alarm threshold that fires on at most a
fixed fraction of no-transition situations, does the detector then fire
on genuine transitions more often than that fraction?} --- a
matched-false-alarm question, followed by a significance test on the
transition hit-rate.

This paper builds the machinery to answer the operational question,
applies it across four independent physical domains and one molecular
modality with one shared protocol, and runs the canonical literature
detector through the identical protocol as a same-segments
head-to-head. The result is a null, and the gap between the two
questions is the finding.

\section{Methods}\label{sec:methods}

\subsection{The detector suite}

One suite of three passive detectors reads a neutral observable bundle
and emits per-window scores:

\begin{itemize}
  \item \textbf{Critical slowing down} --- the larger of the robust
    (median/MAD) z-scores of the windowed variance and lag-one
    autocorrelation of the observable, against a leading baseline.
  \item \textbf{Rising synchronisation} --- the robust z-score of the
    Kuramoto order parameter of a population of phase oscillators.
  \item \textbf{Ordinal-transition entropy} --- the robust z-score of
    the (negated) ordinal-pattern transition entropy of the phase
    field.
  \item \textbf{Weighted fusion} --- the weighted mean of the three
    members.
\end{itemize}

The single-series domain (palaeoclimate) carries one scalar observable,
so only critical slowing down applies; the multi-node domains (EEG,
ECG, grid) carry a population of oscillators and use all four.

\subsection{The matched-false-alarm harness}

For each domain, a per-domain adapter turns the raw recording into the
neutral observable bundle. Each transition is scored on a fixed
pre-onset segment ending at the annotated onset, so every window is
pre-onset and any alarm is a genuine lead. A no-transition null --- a
stable, transition-free stretch --- is cut into non-overlapping trials
of the same length. Each detector's alarm threshold is set continuously
to the tightest value holding the trial false-alarm rate at or below a
target of 10\,\% (the quantile of the null alarm scores, with no grid
ceiling, so a variance-heavy detector is never silently clipped). A
transition is \emph{led} when the detector alarms within its pre-onset
segment. The achieved false-alarm rate is recorded alongside every
threshold.

\subsection{The permutation significance test}

A detection count leaves one question open: does the detector lead
transitions more often than the matched false-alarm rate explains? We
answer it with a label-permutation (exchangeability) test. Each segment
--- transition or null --- alarms or not at the calibrated threshold.
Under the null hypothesis that transition segments are indistinguishable
from nulls, the assignment of the ``transition'' label over the pooled
segments is exchangeable. We draw 10\,000 random transition-sized
subsets of the pooled alarm outcomes (fixed seed, so the p-value is
byte-reproducible), build the null distribution of the lead count, and
report the one-sided p-value with an add-one correction.

\subsection{The competitor head-to-head}

The canonical literature detector is Dakos et al.\ 2008: the Kendall
rank correlation $\tau$ of the windowed lag-one autocorrelation against
time --- a rising-AR(1) trend. We implement it as a per-segment score,
calibrate a matched-false-alarm $\tau$ threshold on the null segments
exactly as above, and run the same permutation test on the resulting
alarms. Because the two detectors read the same one-dimensional signal
(the detrended proxy residual for palaeoclimate, the cross-node
Kuramoto order parameter for the multi-node domains), the head-to-head
is a same-segments, same-budget, same-test comparison in which only the
detector differs.

\subsection{Corpora (citation-only)}

\begin{table}[h!]
\centering
\small
\setlength{\tabcolsep}{4pt}
\begin{tabular}{llcl}
\toprule
Domain & Corpus & Transitions & Null \\
\midrule
Brain (scalp EEG) & CHB-MIT, chb01 (Shoeb 2009, PhysioNet) & 6 seizures & interictal records \\
Heart (ECG) & MIT-BIH AFDB (Moody \& Mark 1983) & 6 AF onsets & sinus stretches \\
Grid (PMU) & PSML 23-bus (Zheng et al.\ 2021) & 12 growing osc. & damped disturbances \\
Palaeoclimate & Dakos et al.\ 2008 records & 6 of 8 evaluated & stable pre-approach \\
\bottomrule
\end{tabular}
\caption{The four physical corpora.}
\end{table}

All raw data are public and cited; none is redistributed. Only derived,
hash-sealed, byte-reproducible evidence records are committed. Each
detector's alarm --- or non-detection --- is sealed into a
content-addressed (canonical-JSON SHA-256) record, so the result,
including every sealed non-detection, is auditable and reproduces
bit-for-bit from a fresh run.

\subsection{The molecular fifth modality (DNB)}

Dynamical-network-biomarker early warning (Chen et al.\ 2012) is a
different modality: few timepoints, many genes, and a cross-sample
statistic computed over the sample population at each timepoint, so the
sliding-window suite does not apply. Its index rises to a peak at the
tipping point. Two published index forms are implemented, each with a
modality-appropriate honest null:

\begin{itemize}
  \item \textbf{Single-cell (Mojtahedi et al.\ 2016).} The transition
    index is the ratio of mean signed gene--gene to mean signed
    cell--cell correlation over the curated panel (which \emph{is} the
    critical module, so no selection step exists). We take the
    per-lineage index trajectory and its bootstrap standard error from
    the published supplement (not re-derived from the confounded raw
    qPCR); the pre-transition rising limb (Days 0, 1, 3) is scored by
    its least-squares slope; the matched-false-alarm null is a
    temporally-shuffled surrogate resampled from the published mean and
    SE; and the three lineages' alarms are tested by the same
    label-permutation core as the physical domains.
  \item \textbf{Bulk (GSE2565 phosgene lung injury).} Here the module
    must be \emph{selected}, which introduces selection freedom: the
    composite index
    $\mathrm{SD_{in}} \cdot \mathrm{PCC_{in}} / \mathrm{PCC_{out}}$ is
    evaluated on a module chosen to peak at the transition, so a rise
    is partly guaranteed by construction. The null therefore re-selects
    the entire module on each surrogate: it shuffles the timepoint
    labels across the arm's samples and reruns module selection from
    scratch, so the surrogate has the same selection freedom as the
    real analysis. The observed rising slope is ranked against these
    selection-controlled surrogates for a one-sided p-value.
\end{itemize}

Both DNB paths seal hash-addressed, byte-reproducible records exactly
as the physical domains do.

\section{Results}\label{sec:results}

\subsection{The suite across four physical domains}\label{sec:fourdomain}

At a matched false-alarm rate (target 10\,\%), detection is sparse
everywhere, and no detector reaches significance in any domain
(Table~\ref{tab:fourdomain}).

\begin{table}[h!]
\centering
\begin{tabular}{llcc}
\toprule
Domain & Best member & Led / N & Permutation $p$ \\
\midrule
Brain (EEG) & critical slowing down & 2 / 6 & 0.215 \\
Heart (ECG) & synchronisation & 2 / 6 & 0.219 \\
Grid (PMU) & critical slowing down & 3 / 12 & 0.195 \\
Palaeoclimate & critical slowing down & 1 / 6 & 0.514 \\
\bottomrule
\end{tabular}
\caption{Best suite member per domain at a matched 10\,\% false alarm.}
\label{tab:fourdomain}
\end{table}

The fusion never leads more transitions than its best single member, so
there is no robust fusion advantage. The observed lead counts are
within what the matched false-alarm rate produces by chance (expected
counts 0.6--1.7 across the domains).

\subsection{Head-to-head: suite vs Dakos AR(1)-Kendall-$\tau$}\label{sec:headtohead}

Running the canonical Dakos detector through the identical protocol on
the identical segments gives Table~\ref{tab:headtohead}.

\begin{table}[h!]
\centering
\begin{tabular}{lcc}
\toprule
Domain & Suite best (led/N, $p$) & Dakos AR(1)-Kendall-$\tau$ (led/N, $p$) \\
\midrule
Palaeoclimate & CSD 1/6, $p = 0.514$ & 0/6, $p = 1.000$ \\
Grid & CSD 3/12, $p = 0.195$ & 1/12, $p = 0.713$ \\
Heart & sync 2/6, $p = 0.219$ & 0/6, $p = 1.000$ \\
Brain (EEG) & CSD 2/6, $p = 0.215$ & 3/6, $p = 0.067$ \\
\bottomrule
\end{tabular}
\caption{Same-segments, same-budget, same-test head-to-head.}
\label{tab:headtohead}
\end{table}

On its own palaeoclimate records --- the data Dakos et al.\ analysed
--- the AR(1)-Kendall-$\tau$ detector, at a matched operating point,
leads zero of six transitions ($p = 1.00$): it does not beat chance,
and it does no better than the suite. The same holds for grid and
heart. The one exception is scalp EEG, where the rising-AR(1) trend is
the strongest signal in the study (3 of 6, $p = 0.067$) --- marginally
short of significance, and better than the suite on that domain.

\subsection{Reading the null}\label{sec:reading}

At a matched operating point with a permutation significance test, no
detector reaches significance at these corpus sizes --- neither the
suite nor the canonical literature detector, in any of the four
domains. Two readings follow. First, the corpora are small: a modestly
skilled detector could remain non-significant at $n = 6$--$12$, so the
null bounds the demonstrated skill rather than establishing absence of
skill (Section~\ref{sec:limitations}). Second, independently of power,
the same detectors on the same data pass the retrospective per-record
test used in the literature and fail the operational one used here; the
choice of test changes the conclusion. The closest approach to
significance --- the AR(1) trend on scalp EEG (3 of 6, $p = 0.067$) ---
is consistent with the seizure-prediction literature's use of
autocorrelation features and warrants a larger EEG corpus.

\subsection{The molecular fifth modality}\label{sec:dnb}

The same conclusion holds when the protocol is carried to a molecular
modality (Table~\ref{tab:dnb}).

\begin{table}[h!]
\centering
\footnotesize
\setlength{\tabcolsep}{4pt}
\begin{tabular}{llll}
\toprule
Case & Test & Result & $p$ \\
\midrule
Single-cell (Mojtahedi 2016) & matched-FA + permutation, 3 lineages & 1 / 3 lineages clear 10\,\% & 0.266 \\
Bulk (GSE2565 phosgene) & selection-controlled surrogate rank & rise does not beat surrogates & 0.385 \\
\bottomrule
\end{tabular}
\caption{The molecular modality under modality-appropriate nulls.}
\label{tab:dnb}
\end{table}

The single-cell transition index rises 2.8-fold toward the erythroid
bifurcation --- a clear effect size --- yet at its published
four-timepoint resolution only the strongest of three lineages clears a
matched false-alarm operating point, and the corpus-level permutation
test is not significant ($p = 0.27$); the binding constraint is the
temporal resolution, not the effect size. In the bulk GSE2565 case, the
phosgene-exposed arm's apparent DNB rise (slope 1.15) sits below the
90th percentile of surrogates that were allowed to re-select the
biomarker module on shuffled time (surrogate-rank $p = 0.385$), and it
barely exceeds the air control's own selected-module rise
($p = 0.414$). On this corpus, granting the null the same selection
freedom as the analysis absorbs most of the apparent signal --- a
result consistent with selection freedom explaining much of the rise,
subject to the caveats of Section~\ref{sec:limitations}. Both molecular
records are hash-sealed and byte-reproducible.

\section{Discussion}\label{sec:discussion}

\textbf{The evaluation gap is the finding.} The difference between
``there is a detectable rising trend in this record's approach''
(retrospective, per-record, vs surrogates) and ``the detector fires on
transitions more than on controls at a fixed false-alarm budget''
(operational, matched, significance-tested) is exactly the gap between
the literature's positive results and this study's null. That gap ---
not any single detector --- is what the four-domain replication and the
Dakos head-to-head expose. A positive early-warning claim intended for
operational use should be required to clear the operational bar; most
published EWS results have cleared only the retrospective one. To be
explicit about scope: at these corpus sizes the nulls bound the
\emph{demonstrated} skill of the detectors tested; they are not
evidence that operational early warning is impossible in these
domains.

\textbf{Selection freedom is a hidden cost of self-selecting
detectors.} The bulk-DNB case adds a specific methodological lesson:
when a detector chooses its own module (or feature set) to peak at the
transition, an honest null must be granted the same choice. A surrogate
that re-selects from scratch on shuffled time absorbs most of the
apparent signal on this corpus --- so a fair null for a self-selecting
detector is not an afterthought but the core of the test.

\textbf{What survives the null.} What this protocol produces even when
detection fails is an auditable evaluation object: a matched
false-alarm operating point, a permutation p-value, and a
content-addressed evidence record for every transition, including every
non-detection. Negative operational results of this kind are currently
rare in the EWS literature, and they are the necessary baseline against
which any claimed operational skill has to be measured. A companion
manuscript (same repository) shows the complementary positive case: a
domain-specific detector that reads a deterministic physical
instability signature --- the growth rate of a power-grid
electromechanical mode --- clears the identical bar decisively, while
the same detector idea on a domain without such a signature does not.

\section{Limitations}\label{sec:limitations}

\begin{itemize}
  \item \textbf{Small corpora, low power.} With 6--12 transitions per
    physical domain --- and 3 lineages (single-cell) or 1 exposed arm
    (bulk) in the molecular modality --- the significance tests have
    limited power; ``at chance'' bounds the demonstrated skill and does
    not prove early warning impossible.
  \item \textbf{One competitor.} Only the Dakos AR(1)-Kendall-$\tau$
    detector was run head-to-head; modern approaches (e.g.\
    learning-based tipping-point predictors) were not tested and could
    behave differently.
  \item \textbf{Single-subject EEG; within-record climate null.} The
    EEG corpus is one subject (chb01); the palaeoclimate null is the
    stable pre-approach interval of each record, which carries that
    record's own variability and so is conservative.
  \item \textbf{Parameterisation.} One reasonable choice of window,
    step, baseline and target false alarm was used per domain; a sweep
    was not performed.
  \item \textbf{DNB caveats.} The single-cell index is taken from the
    published supplement (its bootstrap SE drives the surrogate), not
    re-derived from the raw qPCR; the bulk case has one exposed arm, so
    it is a single-transition surrogate test, not a corpus, and its
    module search is a greedy reimplementation of the DNB selection
    rather than the authors' exact procedure. The molecular conclusion
    is bounded accordingly: no significance at this resolution, and no
    surrogate-beating rise under a selection-controlled null --- not a
    refutation of the DNB method.
\end{itemize}

\section{Reproduction}\label{sec:reproduction}

Every result regenerates deterministically (no randomness in the
pipelines; the permutation test uses a fixed seed). The pipelines are
pure Python (NumPy/SciPy stack; pinned lock files ship in the
repository); a full re-run of the four physical domains and both DNB
cases completes in minutes on a workstation once the cited corpora are
downloaded. With each raw corpus in a directory:

\begin{verbatim}
python bench/early_warning_leadtime_eeg.py     DATA OUT   # brain
python bench/early_warning_leadtime_cardiac.py DATA OUT   # heart
python bench/early_warning_leadtime_grid.py    DATA OUT   # grid
python bench/early_warning_leadtime_climate.py DATA OUT   # palaeoclimate
python bench/head_to_head_ar1_kendall.py OUT \
    --climate-dir C --grid-dir G --cardiac-dir H --eeg-dir E
python -m bench.early_warning_dnb          OUT   # single-cell DNB
python -m bench.early_warning_dnb_bulk     DATA OUT   # bulk GSE2565 DNB
\end{verbatim}

The single-cell DNB case reads only the embedded published summary, so
it needs no external data. The sealed evidence and aggregate results
are committed under \code{examples/real\_data/} in the source
repository (\url{https://github.com/anulum/scpn-phase-orchestrator}). A
fresh run reproduces every \code{content\_hash} bit-for-bit; committed
integrity tests recompute each sealed hash from the payload alone.

\section{Data availability}\label{sec:data}

Raw corpora are public and cited, not redistributed:

\begin{itemize}
  \item CHB-MIT Scalp EEG Database ---
    \url{https://physionet.org/content/chbmit/}
  \item MIT-BIH Atrial Fibrillation Database ---
    \url{https://physionet.org/content/afdb/}
  \item PSML power-system dataset (Zheng et al.\ 2021) --- the 23-bus
    millisecond-level PMU measurements.
  \item Dakos et al.\ 2008 palaeoclimate records ---
    \url{https://github.com/earlywarningtoolbox/datasets}
  \item Mojtahedi et al.\ 2016 single-cell transition index ---
    Table~S2 of the paper (index + bootstrap SE per lineage per day).
  \item GSE2565 phosgene lung-injury expression ---
    \url{https://www.ncbi.nlm.nih.gov/geo/query/acc.cgi?acc=GSE2565}
\end{itemize}

\begin{thebibliography}{9}

\bibitem{scheffer2009} Scheffer M, Bascompte J, Brock WA, et al.
Early-warning signals for critical transitions. \emph{Nature} 461:53
(2009).

\bibitem{dakos2008} Dakos V, Scheffer M, van Nes EH, Brovkin V,
Petoukhov V, Held H. Slowing down as an early warning signal for abrupt
climate change. \emph{PNAS} 105(38):14308 (2008).

\bibitem{dakos2012} Dakos V, Carpenter SR, Brock WA, et al. Methods for
detecting early warning signals of critical transitions in time series.
\emph{PLoS ONE} 7(7):e41010 (2012).

\bibitem{shoeb2009} Shoeb A. Application of machine learning to
epileptic seizure onset detection and treatment. PhD thesis, MIT
(2009); CHB-MIT via PhysioNet.

\bibitem{moody1983} Moody GB, Mark RG. A new method for detecting
atrial fibrillation using R-R intervals. \emph{Computers in Cardiology}
10:227 (1983).

\bibitem{goldberger2000} Goldberger AL, Amaral LAN, Glass L, et al.
PhysioBank, PhysioToolkit, and PhysioNet. \emph{Circulation}
101(23):e215 (2000).

\bibitem{zheng2021} Zheng X, et al. A multi-scale time-series dataset
with benchmark for machine learning in electricity (PSML) (2021).

\bibitem{chen2012} Chen L, Liu R, Liu Z-P, Li M, Aihara K. Detecting
early-warning signals for sudden deterioration of complex diseases by
dynamical network biomarkers. \emph{Scientific Reports} 2:342 (2012).

\bibitem{mojtahedi2016} Mojtahedi M, Skupin A, Zhou J, et al. Cell fate
decision as high-dimensional critical state transition. \emph{PLoS
Biology} 14(12):e2000640 (2016).

\bibitem{sciuto2005} Sciuto AM, Phillips CS, Orzolek LD, Hensley JL,
Chang L-Y, Nadadur SS. Genomic analysis of murine pulmonary tissue
following carbonyl chloride inhalation. \emph{Chemical Research in
Toxicology} 18(11):1654 (2005); GSE2565.

\bibitem{mormann2007} Mormann F, Andrzejak RG, Elger CE, Lehnertz K.
Seizure prediction: the long and winding road. \emph{Brain} 130(2):314
(2007).

\end{thebibliography}

\end{document}
